\documentclass[12pt,a4paper]{article}

% --- Pacchetti di Formattazione ---
\usepackage[utf8]{inputenc}
\usepackage[italian]{babel}
\usepackage[T1]{fontenc}
\usepackage{geometry}
\geometry{margin=2.5cm}
\usepackage{graphicx}
\usepackage{booktabs}
\usepackage{xcolor}
\usepackage{titlesec}
\usepackage{hyperref}
\usepackage{amsmath}
\usepackage{eurosym} 
\usepackage{setspace}
\usepackage{tabularx}

% --- Definizione Colori Branding LIMEN ---
\definecolor{LimenDeep}{HTML}{0D0D1A}     % Blu Notte
\definecolor{LimenAccent}{HTML}{A78BFA}   % Viola Lavanda
\definecolor{LimenSoft}{HTML}{F1F5F9}     % Grigio Ghiaccio

% --- Configurazione Stili ---
\titleformat{\section}{\large\bfseries\color{LimenDeep}}{\thesection}{1em}{}[\titlerule]
\titleformat{\subsection}{\normalsize\bfseries\color{LimenDeep}}{\thesubsection}{1em}{}
\onehalfspacing

\begin{document}

% =========================================================
% COPERTINA PERSONALIZZATA
% =========================================================
\begin{titlepage}
    \centering
    \vspace*{1cm}
    
    \rule{\textwidth}{1.6pt}\vspace*{-\baselineskip}\vspace*{2pt}
    \rule{\textwidth}{0.4pt}\\[\baselineskip]
    
    {\Huge \textbf{\color{LimenDeep} L I M E N}} \\
    \vspace{0.5cm}
    {\Large \textit{Beyond the Threshold of Anxiety}} \\
    
    \rule{\textwidth}{0.4pt}\vspace*{-\baselineskip}\vspace*{2pt}
    \rule{\textwidth}{1.6pt}\\[2cm]
    
    {\Large \textbf{SPECIFICHE TECNICHE E ARCHITETTURA}} \\
    \vspace{0.5cm}
    {\large Bracciale Bionico per la Gestione Predittiva dell'Ansia} \\
    \vspace{0.3cm}
    {\normalsize \textbf{Target:} Edge AI, Wearable Technology, Biofeedback} \\
    
    \vfill

    % SEZIONE TEAM 
    \begin{minipage}{0.8\textwidth}
        \centering
        \textbf{\large IL TEAM DI PROGETTO} \\
        \vspace{0.5cm}
            \textbf{Alberto Bortoletto} \\
            \textbf{Daniele Luconi} \\
            \textbf{Eduardo Pane} \\
            \textbf{Fabrizio Scabbia} \\
            \textbf{Michele Stevanin} \\
            \textbf{Sara Zanella} \\            
            \textbf{Sofia Parsadanova} \\
    \end{minipage}

    \vfill
    {\large \today}
\end{titlepage}

\newpage
\tableofcontents
\newpage

% =========================================================
% CORPO DEL DOCUMENTO
% =========================================================

\section{Architettura Computazionale (Edge AI)}
Per garantire risposte in tempo reale senza dipendere da una connessione smartphone o cloud, LIMEN elabora i dati localmente sul polso utilizzando l'approccio \textbf{TinyML}.

\begin{itemize}
    \item \textbf{Microcontrollore (SoC):} Architettura basata su ESP32-S3 (o equivalente Nordic nRF5340). Scelto per il supporto alle istruzioni vettoriali (accelerazione reti neurali), basso consumo energetico e connettività BLE integrata.
    \item \textbf{Modello Predittivo:} Rete Neurale LSTM (Long Short-Term Memory) a 2 layer. Addestrata per identificare i pattern pre-crisi combinando 5 feature biometriche su finestre scorrevoli di 60 secondi.
    \item \textbf{Compilazione e Ottimizzazione:} Il modello PyTorch nativo viene esportato in ONNX e quantizzato a 8-bit tramite TensorFlow Lite per Microcontrollori (TFLite Micro). Questo riduce il peso del modello a meno di 150 KB, permettendo l'inferenza in pochi millisecondi direttamente sulla RAM del SoC.
\end{itemize}

\section{Layout Hardware}
Il bracciale è ingegnerizzato su due livelli funzionali: la rilevazione (lato pelle) e l'elaborazione/feedback (lato esterno).

\begin{table}[h!]
    \centering
    \small
    \caption{Componentistica Hardware e Posizionamento}
    \vspace{0.3cm}
    \begin{tabularx}{\textwidth}{@{}l l l X@{}}
        \toprule
        \textbf{Componente} & \textbf{Tipologia/Modello} & \textbf{Posizionamento} & \textbf{Funzione Primaria} \\ \midrule
        \textbf{PPG} & MAX30102 (Ottico Verde/IR) & Inferiore, centro & Rilevazione HR e HRV \\
        \textbf{EDA / GSR} & Piastrine in Ag/AgCl & Inferiore, lati PPG & Conduttanza Cutanea \\
        \textbf{Termometro} & MAX30205 (Contatto) & Inferiore, adiacente EDA & Monitoraggio vasocostrizione \\
        \textbf{IMU} & MPU6050 (6-assi) & Scheda Madre (PCB) & Filtraggio artefatti e respiro \\
        \textbf{Motori Aptici} & 3x LRA + Driver DRV2605L & Superiore (12, 4, 8) & Feedback tattile e grounding \\
        \textbf{Alimentazione} & Batteria LiPo 250 mAh curva & Superiore, sopra PCB & Autonomia giornaliera \\
        \bottomrule
    \end{tabularx}
\end{table}

\section{Scelta dei Materiali ed Ergonomia}
Il design di LIMEN privilegia il comfort prolungato e la precisione del segnale biometrico, evitando l'estetica stigmatizzante dei dispositivi medici tradizionali.

\begin{itemize}
    \item \textbf{Scocca Elettronica (Housing):} Realizzata in Policarbonato (PC) leggero o ABS. Garantisce rigidità strutturale per proteggere il PCB e la batteria, mantenendo un peso complessivo inferiore ai 40 grammi.
    \item \textbf{Rivestimento Interno (Skin-contact):} Silicone Medicale Ipoallergenico o TPU (Poliuretano Termoplastico). Isola la scheda madre dal sudore e si adatta all'anatomia del polso, lasciando scoperte solo le lenti ottiche e le piastrine EDA.
    \item \textbf{Cinturino:} Nylon intrecciato elastico con chiusura a micro-regolazione (es. velcro o magnetica). Assicura una pressione radiale costante, requisito fondamentale per non perdere l'aderenza del sensore PPG durante i movimenti.
\end{itemize}

\section{Logica Algoritmica: Il Loop Adattivo Aptico}
Il sistema non si limita a vibrare, ma valuta la risposta fisiologica ogni 40 secondi, modificando la strategia per aggirare l'assuefazione dei meccanocettori (corpuscoli di Pacini).

\subsection{Caso A: De-escalation (Il corpo risponde allo stimolo)}
\begin{itemize}
    \item \textbf{Trigger:} Il SoC rileva l'innalzamento dell'HRV e la stabilizzazione dell'EDA.
    \item \textbf{Azione:} Il pattern aptico rallenta gradualmente (da 0.3 Hz a 0.15 Hz) e l'intensità sfuma fino a spegnersi (Fade-out).
\end{itemize}

\subsection{Caso B: Shift Strategico (L'ansia persiste)}
Se i parametri biometrici non migliorano, l'algoritmo scala l'intensità dell'intervento attraverso 4 fasi distinte:
\begin{enumerate}
    \item \textbf{Fase 1 (0-40s) -- Entrainment Respiratorio:} Sinusoide globale e lenta per sincronizzare meccanicamente il respiro.
    \item \textbf{Fase 2 (40-80s) -- Pressione Crescente:} Rampa a dente di sega per spostare l'attenzione sul polso (Grounding somatico).
    \item \textbf{Fase 3 (80-120s) -- Stimolazione Spaziale Alternata:} I 3 motori LRA si attivano in rotazione per ingannare l'adattamento neurale.
    \item \textbf{Fase 4 (>120s) -- Pattern Interrupt SOS:} Impulsi netti stile codice Morse che costringono il cervello a un'elaborazione cognitiva, distraendo forzatamente l'amigdala.
\end{enumerate}

\end{document}